% p238_1; source PDF page 238; manually verified editable table.
\begin{tikzpicture}[x=1cm,y=-1cm,inner sep=0pt,outer sep=0pt]
\path[use as bounding box] (0,0) rectangle (16.000,9.350);
\draw[line width=.4pt] (0.000,0.000) rectangle (2.000,0.850);
\node[align=center,text width=1.740cm,font={\linespread{1}\fontsize{10}{13.50}\selectfont\bfseries}] at (1.000,0.425) {行为类型};
\draw[line width=.4pt] (2.000,0.000) rectangle (3.800,0.850);
\node[align=center,text width=1.540cm,font={\linespread{1}\fontsize{10}{13.50}\selectfont\bfseries}] at (2.900,0.425) {主观};
\draw[line width=.4pt] (3.800,0.000) rectangle (7.700,0.850);
\node[align=center,text width=3.640cm,font={\linespread{1}\fontsize{10}{13.50}\selectfont\bfseries}] at (5.750,0.425) {结论};
\draw[line width=.4pt] (7.700,0.000) rectangle (10.400,0.850);
\node[align=center,text width=2.440cm,font={\linespread{1}\fontsize{10}{13.50}\selectfont\bfseries}] at (9.050,0.425) {罪数};
\draw[line width=.4pt] (10.400,0.000) rectangle (16.000,0.850);
\node[align=center,text width=5.340cm,font={\linespread{1}\fontsize{10}{13.50}\selectfont\bfseries}] at (13.200,0.425) {范例};
\draw[line width=.4pt] (0.000,0.850) rectangle (2.000,9.350);
\node[align=left,text width=1.740cm,font={\linespread{1}\fontsize{10}{13.50}\selectfont}] at (1.000,5.100) {绑架罪的实行行为（实力控制行为）};
\draw[line width=.4pt] (2.000,0.850) rectangle (3.800,4.550);
\node[align=left,text width=1.540cm,font={\linespread{1}\fontsize{10}{13.50}\selectfont}] at (2.900,2.700) {对重伤结果持故意心理};
\draw[line width=.4pt] (3.800,0.850) rectangle (7.700,4.550);
\node[align=left,text width=3.640cm,font={\linespread{1}\fontsize{10}{13.50}\selectfont}] at (5.750,2.700) {绑架罪+故意伤害罪（重伤）=绑架罪，加重处罚。此时实际是结果加重犯，也即“绑架罪（故意）致人重伤”};
\draw[line width=.4pt] (7.700,0.850) rectangle (10.400,4.550);
\node[align=left,text width=2.440cm,font={\linespread{1}\fontsize{10}{13.50}\selectfont}] at (9.050,2.700) {整体是一个行为，只定绑架罪，加重处罚};
\draw[line width=.4pt] (10.400,0.850) rectangle (16.000,4.550);
\node[align=left,text width=5.340cm,font={\linespread{1}\fontsize{10}{13.50}\selectfont}] at (13.200,2.700) {甲为绑架乙，欲实力控制乙，乙反抗，甲将乙打成重伤，制服了乙，向乙的家人勒索赎金。对甲定绑架罪，加重处罚};
\draw[line width=.4pt] (2.000,4.550) rectangle (3.800,9.350);
\node[align=left,text width=1.540cm,font={\linespread{1}\fontsize{10}{13.50}\selectfont}] at (2.900,6.950) {对死亡结果持过失心理};
\draw[line width=.4pt] (3.800,4.550) rectangle (7.700,9.350);
\node[align=left,text width=3.640cm,font={\linespread{1}\fontsize{10}{13.50}\selectfont}] at (5.750,6.950) {绑架罪和过失致人死亡罪，想象竞合，择一重罪论处};
\draw[line width=.4pt] (7.700,4.550) rectangle (10.400,9.350);
\node[align=left,text width=2.440cm,font={\linespread{1}\fontsize{10}{13.50}\selectfont}] at (9.050,6.950) {一个行为，择一重罪论处，定绑架罪};
\draw[line width=.4pt] (10.400,4.550) rectangle (16.000,9.350);
\node[align=left,text width=5.340cm,font={\linespread{1}\fontsize{10}{13.50}\selectfont}] at (13.200,6.950) {甲用绳子绑住乙，向乙的家人要钱，未要到钱，回来发现乙死了，原来甲不慎将绳子绑得太紧，乙窒息死亡。甲的一个实力控制行为同时触犯绑架罪与过失致人死亡罪，想象竞合，择一重罪论处，定绑架罪};
\end{tikzpicture}
